\documentclass[11pt,a4paper]{article}

\usepackage[margin=1in]{geometry}
\usepackage{booktabs}
\usepackage{graphicx}
\usepackage{hyperref}
\usepackage{xcolor}
\usepackage{amsmath}
\usepackage{multirow}
\usepackage{caption}
\usepackage{enumitem}
\usepackage{float}
\usepackage{url}

\graphicspath{{Figures/}}

\title{Experimental Setup and Analysis:\\
Anomaly Detection in Named Data Networking\\
via Dual Simulation and Emulation}
\author{}
\date{}

\begin{document}
\maketitle

% ============================================================
\section{Introduction}

Named Data Networking (NDN)~\cite{ndn-arch} replaces IP's host-centric addressing with content-centric forwarding: packets carry named content rather than addresses, routers cache responses in a Content Store (CS), the Pending Interest Table (PIT) tracks outstanding requests, and the Forwarding Information Base (FIB) routes interests toward producers. Two attacks exploit this stateful forwarding:

\begin{itemize}[noitemsep]
  \item \textbf{Interest Flooding Attack (IFA):} a compromised consumer floods interests for unresolvable content. The fake interests occupy PIT entries for their full lifetime, saturating routers and blocking legitimate forwarding~\cite{ifa-survey}.
  \item \textbf{Content Poisoning / Cache Pollution (CP):} an attacker requests many unique, rarely-used names. Because the content is served by legitimate producers, it enters router Content Stores through normal caching and continuously evicts popular content~\cite{cp-attack}.
\end{itemize}

We evaluate one \emph{unsupervised} Isolation Forest detector for both attacks on two complementary platforms. \textbf{ndnSIM} (Part~I) is a controlled simulator with deterministic ground truth, used to evaluate both IFA and CP and, via an attacker-rate sweep, to measure each attack's \emph{detection floor}. \textbf{miniNDN} (Part~II) runs real NFD processes and is used as a brief cross-check for IFA. The same detector and the same three topologies are used on both platforms, so agreement across them strengthens validity.

The central result is not merely that the detector works at a high attack rate (it does, at 100\% attacker detection), but that a rate sweep against \emph{live} background traffic exposes a detection floor for each attack, and that the two floors sit an order of magnitude apart. That gap, and the mechanism behind it, is the contribution this document builds toward.

% ============================================================
\part*{Part I: ndnSIM Simulation}

% ============================================================
\section{Simulation Setup}

\paragraph{Environment.} Simulations run in a Docker/Ubuntu~20.04 container using ns-3.35 with ndnSIM~2.x~\cite{ndnsim-repo} (C++17). Scenarios live in \texttt{ndn-research/} and are compiled via Waf; the custom non-responding producer used for IFA lives in \texttt{extensions/}. Two source-level fixes to the patched ndnSIM tree were required: the \texttt{CsTracer} needed the full \texttt{ndn-common.hpp} inclusion to connect to the \texttt{afterCsHit}/\texttt{afterCsMiss} signals (otherwise it reports all zeros), and \texttt{Strategy::sendData()} needed an explicit \texttt{return true;} on its no-PIT-token path (otherwise it crashes on the first CS hit). Both must be re-applied on a fresh checkout, as that subtree is not version-controlled here.

\paragraph{Tracers.} Each run emits four files at a 1\,s interval: \texttt{L3RateTracer} (per-face interest/data/NACK/satisfied/timed-out counts), \texttt{CsTracer} (CS hits/misses), \texttt{AppDelayTracer} (per-packet delay and retransmissions), and a per-second ground-truth label (\texttt{normal}/\texttt{ifa}/\texttt{cp}). Filenames encode the scenario, attacker rate, and seed (Section~\ref{sec:ndnsim-sweep}), so both sweep dimensions are recoverable from the name alone.

% ============================================================
\section{Topologies}

Three topologies are shared with miniNDN (enabling direct comparison). All point-to-point links are 10\,Mbps / 10\,ms unless noted, and every node has a 100-entry Content Store. \textbf{Tree-12}: two producers feed a root router that branches through three routers to six consumers (\texttt{c1}--\texttt{c6}); all consumer traffic crosses the backbone, making it susceptible to network-wide flooding. \textbf{DFN-12}: the same node set with two extra 15\,ms cross-links, forming a denser mesh with redundant paths. \textbf{Dumbbell-10}: three consumers and two producers funnel through a single bottleneck (\texttt{c1}--\texttt{c3} to routers to bottleneck to routers to \texttt{p1},\texttt{p2}), concentrating all traffic and producing the strongest detection signals.

% ============================================================
\section{Traffic Scenarios}

Every scenario runs for 600\,s; attack scenarios inject at $t=300$\,s (balanced 300\,s normal / 300\,s attack), and the detector is fit on the first 200\,s of normal traffic. Routing uses \texttt{GlobalRoutingHelper}.

\paragraph{Legitimate background.} All three scenario types use an identical background so the only difference between a normal and an attack run is the attack itself. Each consumer runs \texttt{ConsumerZipfMandelbrot} at 10\,interests/s under \texttt{/ndn} with a catalog of $N=300$ ($q=0.0$, $s=0.7$); both producers serve 1024-byte payloads. The catalog size is deliberate. A small catalog is exhausted within $\approx 30$\,s at 10\,Hz, after which the consumers fall silent and the background dies, so a model trained on such a run learns ``normal $=$ idle network'' and later measures detection against silence. At $N=300$ the consumers emit for the full 600\,s (verified active at $t=400$\,s, inside the attack window) while the cache stays useful and the simulation tractable. Figure~\ref{fig:traffic-health} confirms the background is alive through the attack window; the impact of getting this wrong is quantified in Section~\ref{sec:ndnsim-floor}.

\begin{figure}[H]
  \centering
  \includegraphics[width=0.7\textwidth]{fig_legit_traffic_health.png}
  \caption{Legitimate consumer interest rate over time for the three topologies, attack window shaded. With the corrected catalog ($N=300$) the background stays flat and alive for the full 600\,s, including through the attack window ($\approx 4.4$\,interests/s per consumer), so every detection number in Section~\ref{sec:ndnsim-floor} is measured against a busy network rather than the silent one that would trivialise both detection and the false-positive rate.}
  \label{fig:traffic-health}
\end{figure}

\paragraph{IFA.} At $t=300$\,s consumer \texttt{c1} launches four concurrent \texttt{ConsumerCbr} streams toward \texttt{/evil} in a fixed $40\!:\!30\!:\!20\!:\!10$ rate ratio (Table~\ref{tab:ifa-streams}), scaled to the chosen total rate (Section~\ref{sec:ndnsim-sweep}). The \texttt{/evil} namespace is served by a custom \texttt{BlackholeProducer} co-located on the producers: it registers the prefix so interests enter the PIT but never returns Data, so each interest stays pending until it times out, producing genuine PIT pressure. Two simpler designs fail: routing under \texttt{/ndn} lets the real producer satisfy the interests (no timeouts), and an unrouted prefix yields an immediate no-route NACK that clears the PIT before any timeout.

\begin{table}[H]
\centering
\caption{IFA attack streams (all from \texttt{c1}); per-stream rates shown at the reference total of 100\,interests/s.}
\label{tab:ifa-streams}
\begin{tabular}{lcc}
\toprule
\textbf{Prefix} & \textbf{Share (rate at 100\,Hz)} & \textbf{InterestLifetime} \\
\midrule
\texttt{/evil/fake/video}      & 40\% (40\,int/s) & 2000\,ms \\
\texttt{/evil/random-attack}   & 30\% (30\,int/s) & 1500\,ms \\
\texttt{/evil/fake/image}      & 20\% (20\,int/s) & 1000\,ms \\
\texttt{/evil/random-attack-2} & 10\% (10\,int/s) & 4000\,ms \\
\bottomrule
\end{tabular}
\end{table}

\paragraph{CP.} At $t=300$\,s two attacker consumers (\texttt{c1},\texttt{c6} in Tree/DFN; \texttt{c1},\texttt{c3} in Dumbbell) each launch 50 \texttt{ConsumerCbr} apps requesting 50 distinct junk names under \texttt{/ndn/junk/}. Because the junk falls under \texttt{/ndn}, the producers answer it, so the junk Data enters router Content Stores through ordinary caching and evicts popular content, raising the legitimate cache-miss rate. The 50-name diversity is fixed across the sweep; only the per-name frequency is scaled to set the attacker's total rate.

% ============================================================
\section{Attacker-Rate Sweep and Replication}
\label{sec:ndnsim-sweep}

The detection-floor analysis requires each attack across a range of intensities, replicated over seeds. Two parameters are exposed on the command line, all else fixed: \texttt{attackRate} (total interests/s, split $40\!:\!30\!:\!20\!:\!10$ for IFA, divided by the 50 junk names for CP, so only volume changes, never name diversity) and \texttt{rngRun} (the seed, via \texttt{RngSeedManager}; without it every run is byte-identical and replication is meaningless). Filenames encode both as \texttt{\{topo\}-\{attack\}-r\{rate\}-run\{seed\}-\{tracer\}}.

The sweep was built in two passes: an initial pass over high/mid rates, and a top-up that added the rates needed to locate each floor (CP filled the $30$--$50$\,Hz gap; IFA was pushed below 10\,Hz). Because the floors are an order of magnitude apart, the two attacks were sampled over different ranges, each dense where its own floor lies (Table~\ref{tab:sweep-matrix}).

\begin{table}[H]
\centering
\caption{ndnSIM attacker-rate sweep. Floor-region rates use five seeds; the two highest rates are single-seed anchors. Rates are per attacker (CP has two attackers).}
\label{tab:sweep-matrix}
\begin{tabular}{lll}
\toprule
\textbf{Scenario} & \textbf{Attacker rates (interests/s)} & \textbf{Seeds} \\
\midrule
Normal & (no attack) & 1--5 \\
\midrule
\multirow{2}{*}{IFA}
  & 1, 2, 4, 6, 8, 10, 12, 15, 20, 30 & 1--5 \\
  & 50, 100 & 1 (anchor) \\
\midrule
\multirow{2}{*}{CP}
  & 10, 12, 15, 20, 30, 35, 40, 45 & 1--5 \\
  & 50, 100 & 1 (anchor) \\
\bottomrule
\end{tabular}
\end{table}

Each configuration runs for all three topologies, giving \textbf{297 runs} (15 normal, 156 IFA, 126 CP), each emitting the four tracer files. A sweep driver records progress and skips configurations whose output already exists, so the campaign is resumable.

% ============================================================
\section{Feature Extraction and Detector}
\label{sec:ndnsim-method}

Raw per-second tracer output is reduced to nine face/cache features per node, all computed as deltas between consecutive 1\,s samples: the three interest/data rates (\texttt{InInterests}, \texttt{OutInterests}, \texttt{InData}) and six ratios (\texttt{cache\_hit\_ratio}, \texttt{satisfaction\_ratio}, \texttt{timeout\_ratio}, \texttt{nack\_ratio}, \texttt{interest\_amp} $=\Delta$Out/$\Delta$In, \texttt{data\_ratio} $=\Delta$InData/$\Delta$In). Each row also carries its \texttt{rate} and \texttt{run} so the floor analysis can group by attacker rate and average over seeds. IFA and CP leave distinct fingerprints (Figure~\ref{fig:eda-timeseries}): IFA collapses \texttt{satisfaction\_ratio} and spikes \texttt{timeout\_ratio} (interests to the black-hole are never satisfied), while CP keeps satisfaction high but drops \texttt{cache\_hit\_ratio} (legitimate content evicted); yet both drive the attacker's interest rate to a large outlier, which is what lets a single direction-agnostic model isolate either attack.

\begin{figure}[H]
  \centering
  \includegraphics[width=0.9\textwidth]{fig_timeseries_tree.png}
  \caption{Mean feature time series on the Tree topology, normal vs.\ IFA vs.\ CP, with the $t=300$\,s attack onset marked. Each attack perturbs a distinct set of counters, motivating the nine-feature representation: IFA drives timeout and satisfaction counters, CP drives cache-hit and volume counters.}
  \label{fig:eda-timeseries}
\end{figure}

The detector is a scikit-learn pipeline of \texttt{StandardScaler} $\to$ \texttt{IsolationForest}~\cite{iforest} ($n\_\text{estimators}=200$, contamination$=0.01$, \texttt{random\_state}$=42$), fit on the nine features using \emph{only} the first 200\,s of the normal scenario, with no attack data and no labels (fully unsupervised). The \texttt{decision\_function} (positive $=$ normal) is mapped to a $[0,100]$ score with $\text{anomaly\_score}=100-\text{norm}$ (higher $=$ more anomalous); an alarm is raised at $\text{anomaly\_score}\geq 70$. Labels are used only for evaluation. A single \emph{global} forest (all topologies pooled) is the primary detector; a per-topology variant is run for comparison.

% ============================================================
\section{ndnSIM Results}

\subsection{Detection at a Representative Rate: Scores, Localisation, and Latency}

At a representative high attacker rate the attacker score jumps far past the alarm (Table~\ref{tab:ndnsim-lift}), reaching the maximum (100) for CP, while benign consumers stay near zero. The onset is immediate: the attacker anomaly score leaves the normal $\pm 1\sigma$ band within about a second of $t=300$\,s (Figure~\ref{fig:score-timeseries}), and every attacker node is detected within $\approx 1$\,s (Figure~\ref{fig:latency}). The global and per-topology forests agree (the latter at 89\% attacker detection).

\begin{table}[H]
\centering
\caption{ndnSIM mean anomaly score before/after attack onset ($t=301$\,s), global model. Benign is the mean over non-attacker consumers.}
\label{tab:ndnsim-lift}
\begin{tabular}{llcccc}
\toprule
\textbf{Topology} & \textbf{Attack} & \textbf{Attacker pre} & \textbf{Attacker post} & \textbf{Benign post} & \textbf{Lift} \\
\midrule
DFN      & IFA & 8.5 & 78.4  & 1.9 & \textbf{69.9} \\
DFN      & CP  & 8.3 & 100.0 & 2.0 & \textbf{91.7} \\
Dumbbell & IFA & 8.7 & 78.4  & 2.7 & \textbf{69.7} \\
Dumbbell & CP  & 8.4 & 100.0 & 3.4 & \textbf{91.6} \\
Tree     & IFA & 8.5 & 78.4  & 1.9 & \textbf{69.9} \\
Tree     & CP  & 8.3 & 100.0 & 2.0 & \textbf{91.7} \\
\midrule
\multicolumn{5}{l}{Normal (all topologies), mean anomaly score} & \textbf{18.4} \\
\bottomrule
\end{tabular}
\end{table}

\begin{figure}[H]
  \centering
  \includegraphics[width=0.85\textwidth]{fig_score_timeseries.png}
  \caption{Isolation Forest anomaly score over time for the attacker nodes against the normal $\pm 1\sigma$ band. For both attacks the attacker crosses the alarm within $\approx 1$\,s of the $t=300$\,s onset and stays above it for the attack window.}
  \label{fig:score-timeseries}
\end{figure}

The detector does not blanket-flag the network. Classifying each node by topological role (via a breadth-first search from the attacker), the unsupervised score forms a clean spatial gradient along the forwarding path: attacker (IFA 78.4 / CP 100.0) $>$ on-path router (75.4 / 88.7) $>$ producer (60.7 / 63.7) $>$ off-path router ($\approx 36$) $>$ benign consumer ($\approx 2$), shown in Figures~\ref{fig:spatial} and~\ref{fig:pernode}. The score follows path membership, not raw hop distance: producers sit several hops away yet alarm because the attack traffic flows to them. This is an implicit attack-path fingerprint without routing-table access. A feature ablation (Figure~\ref{fig:ablation}) confirms the rate/count features carry the signal: a counts-only detector retains near-full attacker detection, while a ratios-only detector degrades.

\begin{figure}[H]
  \centering
  \includegraphics[width=0.85\textwidth]{fig_spatial_localization.png}
  \caption{Mean anomaly score by node role, and score vs.\ hop distance from the attacker. The score traces the attack's forwarding path (attacker $>$ on-path routers $>$ producers $>$ off-path routers $>$ benign consumers) rather than decaying monotonically with hop count.}
  \label{fig:spatial}
\end{figure}

\begin{figure}[H]
  \centering
  \includegraphics[width=0.85\textwidth]{fig_pernode_detection_comparison.png}
  \caption{Per-node detection rate, global vs.\ per-topology model. Attacker and on-path routers are detected at high rates while benign and off-path nodes stay near zero; the global and per-topology models track each other closely.}
  \label{fig:pernode}
\end{figure}

\begin{figure}[H]
  \centering
  \begin{minipage}{0.49\textwidth}
    \centering
    \includegraphics[width=\linewidth]{fig_feature_ablation.png}
    \captionof{figure}{Attacker detection and normal FPR for full, counts-only, and ratios-only feature sets. The rate/count features carry most of the signal.}
    \label{fig:ablation}
  \end{minipage}\hfill
  \begin{minipage}{0.49\textwidth}
    \centering
    \includegraphics[width=\linewidth]{fig_detection_latency.png}
    \captionof{figure}{Detection latency by node role. Attacker nodes are detected within $\approx 1$\,s of onset; on-path nodes follow shortly after.}
    \label{fig:latency}
  \end{minipage}
\end{figure}

\subsection{Detection Floor: How Low Can the Attacker Go?}
\label{sec:ndnsim-floor}

The scores above are measured at a single high rate; against the original $N=100$ background they also give a trivial $0\%$ false-positive rate, because the legitimate consumers are silent by the attack window. To characterise detection \emph{against live traffic} we use the $N=300$ background (Figure~\ref{fig:traffic-health}) and the sweep of Table~\ref{tab:sweep-matrix}, then read off each attack's detection floor: the rate at which attacker-node detection falls below 90\% (Table~\ref{tab:ndnsim-floor}).

\begin{table}[H]
\centering
\caption{ndnSIM attacker-node detection (\%) vs.\ rate, pooled over topologies and seeds, against sustained background traffic. ``--'' $=$ rate not sampled for that attack.}
\label{tab:ndnsim-floor}
\begin{tabular}{lcccccccccc}
\toprule
\textbf{Rate (Hz)} & 1 & 2 & 10 & 20 & 30 & 35 & 40 & 45 & 50 & 100 \\
\midrule
\textbf{CP}  & -- & -- & 21 & 42 & 60 & 71 & 83 & 92 & 100 & 100 \\
\textbf{IFA} & 60 & 99.9 & 100 & 100 & 100 & -- & -- & -- & 100 & 100 \\
\bottomrule
\end{tabular}
\end{table}

The two floors are an order of magnitude apart (Figure~\ref{fig:ndnsim-floor}). \textbf{CP degrades gracefully}, from 100\% above 50\,Hz to $\approx 21\%$ at 10\,Hz, crossing 90\% near $\mathbf{44}$\,\textbf{Hz}, because its only signature is interest \emph{volume}, so it sinks into the legitimate baseline as its rate approaches it. \textbf{IFA holds 100\% down to $\approx 4$\,Hz and only breaks near 1\,Hz} (floor $\approx \mathbf{1.8}$\,\textbf{Hz}), because its signature is \emph{qualitative} (timeouts and NACKs to the black-hole) and survives at $\approx 1\times$ the legitimate per-consumer volume, versus $\approx 10\times$ for CP (Figure~\ref{fig:ndnsim-mechanism}). Detection against sustained traffic carries a realistic $1$--$2\%$ false-positive rate, not the trivial $0\%$ of the idle case (Figure~\ref{fig:fpr}), and the floors are near-identical across topologies (Figure~\ref{fig:heatmap}), a property of the attack type and detector, not the topology. This corrects the earlier ``100\% everywhere'' reading, which reflected catalog exhaustion rather than detector invincibility.

\begin{figure}[H]
  \centering
  \includegraphics[width=0.74\textwidth]{fig_floor_headline.png}
  \caption{ndnSIM detection floor against live traffic. CP (volumetric) has a high floor ($\approx 44$\,Hz); IFA (qualitative timeout/NACK signature) has a floor an order of magnitude lower ($\approx 1.8$\,Hz). Shaded bands are $\pm1$\,SD over topologies and seeds; dashed line marks 90\% detection.}
  \label{fig:ndnsim-floor}
\end{figure}

\begin{figure}[H]
  \centering
  \includegraphics[width=0.66\textwidth]{fig_mechanism_contrast.png}
  \caption{Detection vs.\ attacker volume relative to the legitimate baseline. IFA is detected barely above the crowd ($\approx 1\times$); CP is not reliably caught until $\approx 10\times$, explaining the floor gap. CP at 10\,Hz is louder ($3.7\times$ baseline) than IFA at 1\,Hz ($1.2\times$) yet detected far worse.}
  \label{fig:ndnsim-mechanism}
\end{figure}

Three views make the floor and its mechanism concrete. Figure~\ref{fig:floor-pertopo} shows the per-topology curves with the per-seed spread, confirming the floor is reproducible and topology-independent. Figure~\ref{fig:score-sep} shows why the floor exists in score space: at low CP rate the attacker score distribution collapses into the benign one (no separation, hence 21\% detection), whereas IFA keeps its separation even at 1\,Hz. Figure~\ref{fig:score-rate} plots the same story as mean attacker score vs.\ rate against the alarm line: the CP score slides through the alarm as the rate drops, while IFA stays pinned near 100 until it nearly stops sending.

\begin{figure}[H]
  \centering
  \includegraphics[width=0.9\textwidth]{fig_detection_floor.png}
  \caption{Detection floor per topology (CP left, IFA right) with individual seeds as translucent points, the 90\% line, and the pooled crossing. The three topology curves are essentially superimposed and the per-seed spread is tight.}
  \label{fig:floor-pertopo}
\end{figure}

\begin{figure}[H]
  \centering
  \includegraphics[width=0.9\textwidth]{fig_score_separation.png}
  \caption{Anomaly-score histograms, attacker vs.\ benign: CP at 10\,Hz (missed) and 50\,Hz (detected), IFA at 1\,Hz (floor) and 8\,Hz (detected), with the alarm line. At low CP rate the attacker distribution merges into the benign one; IFA retains separation down to its floor.}
  \label{fig:score-sep}
\end{figure}

\begin{figure}[H]
  \centering
  \begin{minipage}{0.49\textwidth}
    \centering
    \includegraphics[width=\linewidth]{fig_score_vs_rate.png}
    \captionof{figure}{Mean attacker anomaly score vs.\ rate with the alarm line and benign 95th-percentile band. CP crosses the alarm near 30--35\,Hz; IFA stays near 100 until $\approx 2$\,Hz.}
    \label{fig:score-rate}
  \end{minipage}\hfill
  \begin{minipage}{0.49\textwidth}
    \centering
    \includegraphics[width=\linewidth]{fig_fpr_vs_rate.png}
    \captionof{figure}{Benign-consumer false-positive rate vs.\ attacker rate, per topology. FPR is a real, rate-independent $\approx 1$--$2\%$ (tree/DFN $\approx 0.8$--$1.0\%$, dumbbell $\approx 2\%$), not the trivial 0\% of the dead-traffic case.}
    \label{fig:fpr}
  \end{minipage}
\end{figure}

The floor is also a property of the attack type rather than the network: the detection heatmap over (rate $\times$ topology) is near-identical row to row (Figure~\ref{fig:heatmap}). Finally, the feature-shift signature (Figure~\ref{fig:feature-sig}) makes the mechanism explicit: CP drives the volume/data counters hard ($\approx +6\sigma$ on \texttt{InInterests}/\texttt{OutInterests}/\texttt{InData}), a pure-volume signature, whereas IFA's defining shift is \texttt{timeout\_ratio} ($+3.7\sigma$) with depressed satisfaction, a qualitative signature that does not depend on volume. Figure~\ref{fig:timeline} illustrates the two regimes on a single run: at a high rate the score jumps over the alarm the moment the attack starts, and at the floor rate it never reliably crosses.

\begin{figure}[H]
  \centering
  \includegraphics[width=0.9\textwidth]{fig_floor_heatmap.png}
  \caption{Detection (\%) over attacker rate $\times$ topology for CP (left) and IFA (right). The three topology rows are near-identical at every rate, so the floor is a property of the attack type and detector, not the topology.}
  \label{fig:heatmap}
\end{figure}

\begin{figure}[H]
  \centering
  \begin{minipage}{0.49\textwidth}
    \centering
    \includegraphics[width=\linewidth]{fig_feature_signature.png}
    \captionof{figure}{Attacker mean feature shift in the attack window (std-devs from normal). CP moves volume and cache counters; IFA moves \texttt{timeout\_ratio} and satisfaction, a qualitative fingerprint.}
    \label{fig:feature-sig}
  \end{minipage}\hfill
  \begin{minipage}{0.49\textwidth}
    \centering
    \includegraphics[width=\linewidth]{fig_score_timeline.png}
    \captionof{figure}{Attacker score over 600\,s (8\,s smoothing), high vs.\ floor rate. At high rate the score clears the alarm at onset; at the floor rate it hovers at the alarm and detection is unreliable.}
    \label{fig:timeline}
  \end{minipage}
\end{figure}

% ============================================================
\part*{Part II: miniNDN Emulation (IFA cross-check)}

% ============================================================
\section{miniNDN Cross-Check}

miniNDN~\cite{minindn} runs real NFD forwarding daemons on Mininet network namespaces with NLSR routing, capturing system-level effects (OS scheduling, NFD buffering, timing jitter) absent from simulation. Per-node NFD management counters (including PIT and CS sizes) are polled at $\approx 10$\,s intervals to JSONL logs; a 10-minute warmup is discarded. The campaign covers \textbf{IFA only} (CP was not collected, a limitation; ndnSIM provides CP coverage). Ten delta-based features are derived per node, and the \emph{same} Isolation Forest configuration is trained on clean normal traffic only, with the raw score mapped to a $[0,100]$ scale (30\% strict / 50\% operational thresholds).

\begin{table}[H]
\centering
\caption{miniNDN IFA: attacker \texttt{c1} detection rate and normal-traffic FPR (global model).}
\label{tab:minindn-detection}
\begin{tabular}{lcccc}
\toprule
\multirow{2}{*}{\textbf{Topology}} &
\multicolumn{2}{c}{\textbf{Attacker \texttt{c1} detection}} &
\multicolumn{2}{c}{\textbf{Normal FPR}} \\
\cmidrule(lr){2-3}\cmidrule(lr){4-5}
& @30\% & @50\% & @30\% & @50\% \\
\midrule
Tree     & 89.1\% & 91.2\% & 0.19\% & 1.08\% \\
Dumbbell & 98.5\% & 98.7\% & 1.00\% & 4.27\% \\
DFN      & 96.7\% & 96.8\% & 0.64\% & 4.10\% \\
\bottomrule
\end{tabular}
\end{table}

Attacker detection exceeds 91\% at the operational threshold across all topologies, at $<4.3\%$ false positives on clean data (Table~\ref{tab:minindn-detection}). Detection again traces the forwarding path: backbone routers and on-path producer \texttt{p1} alarm at near-attacker rates (88--90.5\% in Tree) while off-path leaf routers stay low. One emulation-specific effect appears in the \textbf{Tree} topology, where benign consumers flag at 51--58\% (versus $<7.3\%$ in DFN/Dumbbell): because the hierarchy forces all traffic through the shared backbone, \texttt{c1}'s flood degrades forwarding for every branch, making IFA a network-wide QoS event rather than a localised one. KS tests confirm unambiguous separation on the attacker node (cache\_hit\_ratio, satisfaction\_ratio, nack\_rate all $D\!>\!0.93$); \texttt{pit\_growth\_rate} is non-discriminating ($D=0.028$) at the coarse 10\,s polling. As in ndnSIM, the global model matches per-topology models, so no topology-specific calibration is needed.

% ============================================================
\section{Key Findings and Limitations}

\begin{enumerate}[noitemsep]
  \item \textbf{One unsupervised Isolation Forest detects both IFA and CP from face-level counters (ndnSIM), with no PIT access.} At high rate it isolates every attacker node (100\% detection, $\approx 1$\,s latency), so no attack-specific detector is needed.
  \item \textbf{The two attacks have detection floors an order of magnitude apart.} Against live traffic, CP crosses 90\% near \textbf{44\,Hz} (purely volumetric) while IFA holds to $\approx 4$\,Hz and breaks near \textbf{1.8\,Hz} (qualitative timeout/NACK signature), at a realistic $1$--$2\%$ FPR.
  \item \textbf{The floor gap is explained mechanistically.} A volumetric attack must out-shout the legitimate baseline ($\approx 10\times$ for CP), so it hides as its rate approaches the crowd; a qualitative attack stays detectable at $\approx 1\times$ the baseline until it nearly stops sending. This is a cleaner, more defensible claim than ``100\% everywhere.''
  \item \textbf{Signatures are topology-invariant (both platforms).} The global model matches per-topology models across hierarchical, mesh, and bottleneck structures without retraining, so a detector trained at one site transfers to another.
  \item \textbf{Detection is spatially localised (both platforms).} The score forms a gradient along the attack's forwarding path, giving an implicit attacker fingerprint without routing-table access.
  \item \textbf{Cross-platform agreement strengthens validity.} miniNDN (real NFD, PIT-based features) independently reproduces the topology-invariant, path-localised IFA detection seen in ndnSIM.
\end{enumerate}

\noindent\textbf{Limitations.} CP was not evaluated in miniNDN (ndnSIM covers it). Victim consumers are not directly detected (their own traffic looks normal). Reported floor values are interpolated crossings and are best quoted as bands (CP 40--45\,Hz, IFA 1--2\,Hz) rather than false-precision decimals; the two highest sweep rates are single-seed anchors. ndnSIM uses idealised links, and miniNDN's 10\,s polling smooths fast transients (PIT growth becomes non-discriminating). Concurrent mixed attacks and a systematic threshold-sensitivity analysis were not performed, though the detector is confirmed stable across seeds and the global/per-topology choice.

% ============================================================
\begin{thebibliography}{9}

\bibitem{ndnsim-repo}
R.~Bhusal, ``Patched ndnSIM for ns-3.35,'' GitHub repository, 2024. [Online]. Available: \url{https://github.com/bhusal-rj/ndnSIM-ns3}

\bibitem{ndn-arch}
L.~Zhang et al., ``Named data networking,'' \textit{ACM SIGCOMM CCR}, vol.~44, no.~3, pp.~66--73, 2014.

\bibitem{minindn}
S.~Mastorakis et al., ``ndnSIM 2.0: A new version of the NDN simulator for NS-3,'' \textit{NDN Technical Report NDN-0028}, 2016.

\bibitem{iforest}
F.~T.~Liu, K.~M.~Ting, and Z.-H.~Zhou, ``Isolation forest,'' in \textit{Proc.\ IEEE ICDM}, 2008, pp.~413--422.

\bibitem{ifa-survey}
A.~Compagno, M.~Conti, P.~Gasti, and G.~Tsudik, ``Poseidon: Mitigating interest flooding DDoS attacks in named data networking,'' in \textit{Proc.\ IEEE LCN}, 2013, pp.~630--638.

\bibitem{cp-attack}
C.~Ghali, G.~Tsudik, and E.~Uzun, ``Needle in a haystack: Mitigating content poisoning in named-data networking,'' in \textit{Proc.\ NDSS SENT Workshop}, 2014.

\end{thebibliography}

\end{document}
